% Draft: §5, once, twice and thrice, for the restructured paper (rewrite
% method agreed 17 Oct 2026; see the evidence repo's DECISIONS.md). Drafted
% 18 Oct from Payne et al. (2007) pp. 584--597, read in full from the PDF, and
% the structure review's plan (reviews/structure-muse-20261017-115410/MUSE.md,
% "Payne et al.'s own case"; budget 800 words in the architecture table).
% The old draft's subsection (fusion-of-functions-revisited.tex, sec:once) was
% written on version B (once heading an NP that functions as modifier in the
% once mayor) and isn't carried over; version A throughout.
% All examples are Payne et al.'s; no new corpus data.
% Revised 18 Oct after the muse review (reviews/muse-sec5-20261018-101536/
% TRIAGE.md): multiplier, n. 40 and the non-standard uses (pp. 596--597) added;
% the PP up to twice answered from n. 31; singulative gaps put as the §4 default
% plus the adjective trade; term of office conceded the restriction to offices
% (pp. 585, 593); every page rechecked against the PDF text.
%
% Labels it relies on: sec:analyses (§2), sec:det-head (§3), sec:mod-head (§4).
% Owes §6 a row (drafts/costs.tex) and a clause in the "second function" row.

\section{\mentionhead{Once}, \mentionhead{twice}, and \mentionhead{thrice}}\label{sec:once}

PHP built their case for fusion on \mention{once}, \mention{twice}, and \mention{thrice}. The uses PHP analyse with fusion are numerical, singulative (\mention{Not once have I seen anything like this}), past (\mention{I was once a teenage werewolf}), and term of office (\mention{the once mayor}); sequential \mention{once}, as in \mention{once you have finished}, is a preposition (pp.~584, 592--594).

Much of their argument concerns category, and it shows that numerical \mention{once} is a determinative (pp.~586--589): it takes determiners, relative clauses, and the modifiers of \mention{one} (\mention{more than once}), and it can be a subject (\mention{Once is enough}). Ordinary Head agrees, so only the second function is in question.

PHP take the form of the word to show that function too. In \mention{someone} and the other compound determinatives, \enquote{the morphological compounding in this case reflects the fusion of two syntactic functions}, and \mention{once} is \enquote{similarly} a compound, of a determinative base \mention{on·} and a noun base \mention{·ce} \enquote*{time} (pp.~587--588). But simple \mention{some} and \mention{few} are fused with no compounding, and PHP's own note suggests \mention{·ce} is a reanalysed genitive ending (p.~588 n.~30), and \textcite[89]{wu2021exocentric} finds their noun base \enquote{far from convincing}. The two bases can as well reflect the word's meaning, a quantifier with a sort, as in the compounds (Section~\ref{sec:det-head}), and both analyses give \mention{once} that meaning.

The first thing fusion is said to explain is why \mention{once} can head a nominal with no determiner as well as with one, which a singular count noun normally can't (\ref{ex:dog}b). Numerical \mention{one} functions as determiner (\mention{one time}) and as modifier after a determiner (\mention{this one time}), so \mention{once} is a fused determiner--head on its own and a fused modifier--head after \mention{this} (p.~588). That's the link between the two uses that Section~\ref{sec:mod-head} examined for determinatives, a default with exceptions both ways. Ordinary Head states it as a permission: \mention{once} may head a nominal with or without a determiner.

The second is the position of an adjective. Since the fused \mention{once} is also the determiner, \ungram{\mention{last once}} is excluded \enquote{for the same reason as} \ungram{\mention{last one time}}, while \mention{this last once} is fine (p.~589). One of their \textit{JL} referees judged \mention{this last once} ungrammatical, and they allow that speakers may vary (p.~586 n.~26). My searches of COCA, NOW, and iWeb found no relevant example of \mention{this last once} or \mention{the last once}, but they found the pattern with \mention{twice}, their (33b) (p.~595), as in \mention{He's gone well \bk{the last \fu{twice}}} (\textit{Racing Post}, 2026, NOW).\footnote{The hits for \mention{once} are the verb \mention{last}, slips for \mention{one} or \mention{ounce}, and adverbial or sequential \mention{once}. \mention{The last twice} is common in NOW and iWeb, mostly in British and Irish racing reports, and absent from COCA; the counts are in the evidence index.} So the adjective after a determiner is attested for \mention{twice}, not \mention{once}.

The gap before a bare \mention{once} is the one Section~\ref{sec:mod-head} weighed for \mention{few}. Fusion derives it from adjacency, and ordinary Head excludes it by its condition on pre-head modifiers (Section~\ref{sec:mod-head}). PHP's further example, \ungram{\mention{\pn{this} once last}} with a silent head after \mention{last} (p.~589), separates fusion from a silent head, not from ordinary Head.

Larger phrases go the same way. For PHP, the DP \mention{more than once} is a fused determiner--head on its own and a fused modifier--head after a determiner (\mention{the more than once which would have been permitted}, pp.~589--591), and the PP \mention{up to twice} is a fused determiner--head (p.~589 n.~31). Ordinary Head gives each of these phrases a single function.

Singulative and past \mention{once} take no determiner and no attributive adjective (p.~592), and both gaps fall under what's already been weighed: the missing determiner as a lexical restriction stated on both analyses, the adjective under the comparison Section~\ref{sec:mod-head} made for \mention{few}.

Term-of-office \mention{once} functions as attributive modifier (\mention{the once mayor}), as does \mention{twice} in \mention{the twice president}, which many speakers find unnatural (pp.~593, 595). Both analyses have to state that the use is largely confined to nouns for offices (\ungram{\mention{my once job as a glass-blower}}, pp.~585, 593). For PHP, \mention{once} here \enquote{will act as a fused modifier-head forming a nominal rather than forming a full NP} (p.~593), so the attribute is a nominal whose head is also its modifier. On ordinary Head, the determinative \mention{once} simply functions as modifier in the nominal \mention{once mayor}, as \mention{two} does in \mention{the two books}. That's one function where fusion has two, and one nominal fewer, though the lexical statement is no shorter.

The rest of their case (multiplier \mention{twice}, \mention{once a week}, and the non-standard uses, pp.~595--597) turns on category and on the word heading a nominal, which the analyses share.

Apart from the adjective restriction weighed in Section~\ref{sec:mod-head}, no property of form in this paradigm varies with the second function rather than with the word, so on the criterion of Section~\ref{sec:analyses} the paradigm fusion was developed on adds no evidence of its own.\footnote{COCA, NOW, and iWeb have no genitive of \mention{once}, \mention{twice}, or \mention{thrice}. Ordinary Head adds them to its list of gaps. Fusion predicts the gap for the uses PHP analyse as Det--Head, unless a compound's nominal base can carry the genitive, as in \mention{everyone's} (Section~\ref{sec:det-head}), and then it would have to set \mention{·ce} apart from \mention{-one}.}
